% ============================================================================
%  비아벨 Artin L-함수에 대한 ξ-다발 프레임워크
%  와 곡률 판별자의 임계선 외 보편성
%  (프로젝트 RDL — 제3논문)
%  결과 #121–#124
%  ── 한국어판 ──
%  컴파일: xelatex artin_master_ko.tex (2회 실행)
% ============================================================================
\documentclass[11pt, a4paper]{article}

% ── 한글 지원 (XeLaTeX + kotex) ────────────────────────────────────────────
\usepackage{kotex}
\usepackage{fontspec}

% ── Layout ─────────────────────────────────────────────────────────────────
\usepackage[top=2.5cm, bottom=2.5cm, left=2.8cm, right=2.8cm]{geometry}
\usepackage{setspace}
\onehalfspacing

% ── Mathematics ─────────────────────────────────────────────────────────────
\usepackage{amsmath, amssymb, amsthm, mathtools}
\usepackage{bm}

% ── Tables ──────────────────────────────────────────────────────────────────
\usepackage{booktabs}
\usepackage{array}
\usepackage{multirow}
\usepackage{colortbl}
\usepackage{longtable}

% ── Colors ──────────────────────────────────────────────────────────────────
\usepackage[dvipsnames,svgnames,x11names]{xcolor}
\definecolor{txblue}{HTML}{2563EB}
\definecolor{chgreen}{HTML}{059669}
\definecolor{rxorange}{HTML}{D97706}
\definecolor{lossred}{HTML}{DC2626}
\definecolor{decodepurple}{HTML}{7C3AED}
\definecolor{secgray}{HTML}{374151}
\definecolor{lightbg}{HTML}{F8FAFC}
\definecolor{accentblue}{HTML}{3B82F6}
\definecolor{zetaviolet}{HTML}{6D28D9}
\definecolor{primegold}{HTML}{B45309}

% ── Cross-references & Links ───────────────────────────────────────────────
\usepackage[colorlinks=true, linkcolor=accentblue, citecolor=chgreen,
            urlcolor=txblue]{hyperref}
\usepackage[capitalise, noabbrev]{cleveref}

% ── Theorem Environments (한국어, 제1·2논문 번호 연속: 정리 1–7 이후) ────
\theoremstyle{plain}
\newtheorem{theorem}{정리}
\setcounter{theorem}{7}               % 다음 정리는 정리 8
\newtheorem{lemma}[theorem]{보조정리}
\newtheorem{proposition}[theorem]{명제}
\newtheorem{corollary}[theorem]{따름정리}
\newtheorem{conjecture}[theorem]{추측}
\newtheorem{observation}[theorem]{관찰}

\theoremstyle{definition}
\newtheorem{definition}{정의}[section]
\newtheorem{property}{성질}[section]
\newtheorem{remark}{비고}[section]
\newtheorem{example}{예제}[section]

% cleveref 한국어 이름 설정
\crefname{theorem}{정리}{정리}
\Crefname{theorem}{정리}{정리}
\crefname{observation}{관찰}{관찰}
\Crefname{observation}{관찰}{관찰}
\crefname{corollary}{따름정리}{따름정리}
\Crefname{corollary}{따름정리}{따름정리}
\crefname{lemma}{보조정리}{보조정리}
\Crefname{lemma}{보조정리}{보조정리}
\crefname{definition}{정의}{정의}
\Crefname{definition}{정의}{정의}
\crefname{remark}{비고}{비고}
\Crefname{비고}{비고}{비고}
\crefname{section}{§}{§}
\Crefname{section}{§}{§}
\crefname{appendix}{부록}{부록}
\Crefname{appendix}{부록}{부록}
\crefname{table}{표}{표}
\Crefname{table}{표}{표}
\crefname{equation}{식}{식}
\Crefname{equation}{식}{식}

% ── Custom Commands ────────────────────────────────────────────────────────
\newcommand{\norm}[1]{\left\|#1\right\|}
\newcommand{\abs}[1]{\left|#1\right|}
\newcommand{\ie}{\textit{i.e.}}
\newcommand{\eg}{\textit{e.g.}}
\newcommand{\sectionline}{\noindent\rule{\textwidth}{0.4pt}\vspace{0.3em}}

% 다발 관련 명령어
\newcommand{\Lam}{\Lambda}             % 완성 L-함수
\newcommand{\kap}{\kappa}             % 다발 곡률
\newcommand{\mono}{\mathcal{M}}       % 모노드로미
\newcommand{\FE}{\mathrm{FE}}         % 함수방정식
\newcommand{\cv}{\mathrm{CV}}         % 변동계수
\newcommand{\DHf}{\mathrm{DH}}         % 다벤포트-하일브론
\newcommand{\EC}{\mathrm{EC}}         % 타원곡선
\newcommand{\RS}{\mathrm{RS}}         % 랭킨-셀베르그
\newcommand{\Ep}{\mathrm{Ep}}         % 엡스타인
\newcommand{\Gal}{\mathrm{Gal}}       % 갈루아군
\newcommand{\Frob}{\mathrm{Frob}}     % Frobenius

% ── Box Style ──────────────────────────────────────────────────────────────
\usepackage{tcolorbox}
\tcbuselibrary{skins, breakable}

% ============================================================================
\begin{document}

% ── 제목 ──────────────────────────────────────────────────────────────────
\begin{center}
  {\LARGE\bfseries 비아벨 Artin $L$-함수에 대한}\\[4pt]
  {\LARGE\bfseries $\xi$-다발 프레임워크}\\[6pt]
  {\large\bfseries 와 곡률 판별자의 임계선 외 보편성}\\[18pt]
  {\large 김영후}\\[4pt]
  {\normalsize 독립 연구자}\\[4pt]
  {\small\today}\\[4pt]
  {\small\textit{($\xi$-다발 시리즈 제3논문; 동반 논문: Kim 2026a,b)}}
\end{center}

\vspace{1.5em}

% ── 초록 ───────────────────────────────────────────────────────────────────
\begin{tcolorbox}[
  colback=lightbg, colframe=zetaviolet!60!black,
  title={\bfseries 초록}, fonttitle=\large,
  arc=3pt, boxrule=0.8pt, toptitle=3pt, bottomtitle=3pt
]
$\xi$-다발 프레임워크 (Kim~2026a,b)를 비아벨 Artin $L$-함수로 확장하고,
곡률 판별자 $c_{1}$의 임계선 외 보편성이 복수의 $L$-함수 족에 걸쳐
성립함을 확립한다.

$S_{3}$의 표준 표현 ($\mathrm{GL}(2)$, $N=23$)과 $S_{4}$의 표준 표현
($\mathrm{GL}(3)$, $N=283$)에 대해 네 가지 $\xi$-다발 성질
--- 함수방정식 일치, 영점 열거, 기울기 $-0.994\pm0.001$과
$-0.993\pm0.003$의 $\kap\delta^{2}$-스케일링 ($t$-방향), 그리고 모노드로미 $\pm\pi$ ---
을 모두 검증한다. 이는 비아벨 갈루아 표현에 대한 $\xi$-다발 프레임워크의
최초 수치 검증이다.

나아가, PARI/GP \texttt{lfundiv}를 사용한 $S_{3}$의 $\sigma$-방향 재분석
(결과~\#207)은 기울기 $2.0000\pm0.0000$ ($R^{2}=1.000000$, 5개 영점,
FE${}=-197$)을 산출하여, 임계선 위 영점에서의 $c_{1}=0$ 보편성을
$\sigma$-방향에서 확인한다.
$S_{5}$ Artin $L$-함수의 새로운 측정 (차수~$4$, $|S_{5}|=120$; 결과~\#209)은
기울기 $1.9999\pm0.0003$ ($R^{2}=1.000000$, 5개 영점, 모노드로미 $5/5=2.0000\pi$,
$\sigma$-유일성 $5/5$ 통과)을 산출하여, 두 번째 아틴 족에서
$\kap\delta^{2}$ 보편성을 확인하고 차수~$4$에 대한 사슬-편향 우려를
직접 반박한다.
두 결과는 모두 아틴 족을 제2논문의 14개 $L$-함수 비교표
(차수 1--6, 모두 기울기~$\approx 2.0$)와 완전히 일치시킨다.
이 이중-방향 구조 --- $t$-방향 기울기~$\approx{-1}$ (개별 해석적 구조)와
$\sigma$-방향 기울기~$=2.0$ (보편적 $c_{1}=0$ 법칙) --- 는
비아벨 Artin $L$-함수에 대한 $\xi$-다발 기하학의 더 풍부한 특성화를 제공한다.

아울러, Davenport--Heilbronn 함수 (mod~5, 임계선 외 영점 $N=4$)에서
확인된 정리~7의 임계선 외 판별자 $c_{1}\cdot\abs{\sigma_{0}-\tfrac{1}{2}} \to 1$을,
디리클레 지표 mod~7과 mod~11로부터 구성한 일반화 DH 함수로 확장하여,
세 개의 독립적인 $L$-함수 족에 걸친 $N=12$개의 임계선 외 영점을 얻는다.
확장된 측정 결과는 $c_{1}\cdot\abs{\sigma_{0}-\tfrac{1}{2}} = 1.136 \pm 0.162$로,
도체 독립적 보편성을 동반한 점근 법칙을 확인한다.

\smallskip
\noindent 여섯 개의 수치 결과 (\#121--\#124, \#207, \#209)가 비아벨 산술로의 확장,
곡률 판별자의 보편성, 두 갈루아 군 ($S_{3}$, $S_{5}$)에 걸친
아틴 족 보편성을 지지한다.
\end{tcolorbox}

\vspace{0.5em}
\noindent\textbf{핵심어:} $\xi$-다발, Artin $L$-함수,
비아벨 갈루아 표현, $S_{3}$, $S_{4}$,
$\kap\delta^{2}$ 곡률, Davenport--Heilbronn 함수,
임계선 외 영점, 곡률 판별자, 모노드로미, 리만 가설

\tableofcontents
\newpage


% ============================================================================
\section{서론}
\label{sec:intro}
% ============================================================================

Kim~(2026a) (이하 \textit{제1논문})에서 도입되고 Kim~(2026b)
(\textit{제2논문})에서 확장된 $\xi$-다발 프레임워크는,
$L$-함수의 영점을 연구하기 위한 기하학적 시각을 제공한다.
핵심 관측량은 완성 $L$-함수 $\Lam(s)$에 대응하는 선다발의
곡률 불변량 $\kap = \abs{\Lam'/\Lam}^{2}$과,
각 영점 주위의 모노드로미 $\mono = \Delta\arg(\xi)$이다.

제1·2논문은 광범위한 $L$-함수 족에 걸쳐 네 가지 보편적 성질 (P1--P4)을 확립하였다:
\begin{itemize}
  \item \textbf{P1} (함수방정식): $\Lam(s) = \varepsilon\,\overline{\Lam}(1-\bar{s})$;
  \item \textbf{P2} (영점 개수): 편각 원리에 의한 영점 열거;
  \item \textbf{P3} ($\kap\delta^{2}$ 스케일링): 임계선 위에서 $\kap\delta^{2} = 1 + O(\delta^{2})$;
  \item \textbf{P4} (모노드로미): 단순 영점마다 임계선을 따라 $\pm\pi$ 기여.
\end{itemize}
이는 리만 $\zeta$-함수, $\mathrm{GL}(1)$--$\mathrm{GL}(5)$ 자기형태 $L$-함수,
타원곡선 $L$-함수 (계수 0--3), 랭킨-셀베르그 $\mathrm{GL}(4)$,
데데킨트 $\zeta$-함수, 엡스타인 $\zeta$-함수 --- 합계 92개의 수치 결과 ---
에 대해 검증되었다.

제2논문은 나아가 \textbf{정리~7}을 증명하였다:
$\Lam(s)=\varepsilon\,\overline{\Lam}(1-\bar{s})$를 만족하는
$L$-함수의 영점 $\rho = \sigma_{0}+it_{0}$에 대해,
\[
  \kap\delta^{2} = 1 + c_{1}\delta + O(\delta^{2}),
  \qquad c_{1} = \mathrm{Re}\!\left(\frac{\Lam''(\rho)}{\Lam'(\rho)}\right),
\]
이며 $c_{1}=0$이면 그리고 그때에만 $\sigma_{0}=\tfrac{1}{2}$이다.
이는 Davenport--Heilbronn 함수 (mod~5)의 네 개의 알려진 임계선 외 영점에서 검증되었다.

\paragraph{본 논문의 목표.}
두 가지 상호 보완적 확장을 추구한다.

\textbf{확장~A: 비아벨 갈루아 표현.}
제1·2논문에서 검사된 모든 $L$-함수는 아벨 지표 또는 자기-쌍대 모듈형식에서 유래한다.
본 논문에서는 $S_{3}$와 $S_{4}$의 \emph{표준 표현}에 부속된 Artin $L$-함수에 대해
네 가지 성질을 검사한다:
$S_{3}$ ($\mathrm{GL}(2)$, 차수~2)와 $S_{4}$ ($\mathrm{GL}(3)$, 차수~3).
이는 $\xi$-다발 프레임워크에 처음으로 적용되는 비아벨 갈루아 표현이다.

\textbf{확장~B: $c_{1}$ 보편성.}
디리클레 지표 mod~7과 mod~11로부터 일반화 Davenport--Heilbronn 함수를 구성하고,
임계선 외 영점을 찾아 $c_{1}\cdot\abs{\sigma_{0}-\tfrac{1}{2}} \to 1$이
세 개의 독립적인 $L$-함수 족 ($N=12$개 영점)에 걸쳐 성립함을 검증한다.


% ============================================================================
\section{프레임워크 요약}
\label{sec:recap}
% ============================================================================

제1·2논문의 필수 정의와 정리를 간략히 재서술한다;
상세한 내용은 원 논문을 참조하라.

\begin{definition}[$\xi$-다발]\label{def:bundle}
$\abs{\varepsilon}=1$인 함수방정식 $\Lam(s) = \varepsilon\,\overline{\Lam}(1-\bar{s})$를
만족하는 완성 $L$-함수 $\Lam(s)$에 대해, \emph{$\xi$-다발}은
임계대 $\mathcal{S} = \{s \in \mathbb{C} : 0<\mathrm{Re}(s)<1\}$ 위의
$\mathrm{U}(1)$-다발로, 단면 $\psi(s) = \Lam(s)/\abs{\Lam(s)}$,
접속 $A = (\Lam'/\Lam)\,ds$, 곡률 밀도 $\kap(s) = \abs{\Lam'/\Lam}^{2}$를
갖는다.
\end{definition}

\noindent\textbf{정리 1--4} (제1논문):
\begin{enumerate}
  \item 접속 반대칭성: $L(1-s) = -L(s)$, 임계선 위에서 $\mathrm{Re}(L)=0$ (유니터리 게이지).
  \item 편각 원리에 의한 위상적 영점 계산.
  \item 곡률 집중: $\kap \sim m^{2}/\abs{s-\rho}^{2}$, 차수-$m$ 영점 근방;
    일반적인 임계선 점에서 곡률은 소멸.
  \item $\kap\delta^{2}$ 판별자 (정리~4): $c_{1}=0 \iff \sigma_{0}=\tfrac{1}{2}$.
\end{enumerate}

\noindent\textbf{정리 5--7} (제2논문):
\begin{enumerate}\setcounter{enumi}{4}
  \item 타원곡선, 데데킨트, 엡스타인 족에 걸친 보편성.
  \item FE-전용 충분성: $\xi$-다발은 함수방정식 구조에만 의존하며 오일러 곱에는 의존하지 않는다.
  \item 임계선 외 판별자: $c_{1} = \mathrm{Re}(\Lam''/\Lam')$, 점근 법칙
    $c_{1}\cdot\abs{\sigma_{0}-\tfrac{1}{2}} \to 1$.
\end{enumerate}

검증 가능한 네 가지 성질 P1--P4가 \cref{tab:properties}에 요약되어 있다.

\begin{table}[ht]
\centering
\caption{네 가지 $\xi$-다발 성질과 검증 기준.}
\label{tab:properties}
\begin{tabular}{@{}clll@{}}
\toprule
\textbf{성질} & \textbf{이름} & \textbf{판정 기준} & \textbf{통과 조건} \\
\midrule
P1 & 함수방정식 & $\abs{\Lam(s)/\Lam(1-\bar{s}) - \varepsilon} < 10^{-6}$ & 오차 $< 10^{-6}$ \\
P2 & 영점 개수 & $Z(t)$ 부호 변환에 의한 영점 & 일치하는 개수 \\
P3 & $\kap\delta^{2}$ 스케일링 & $\log\kap$ 대 $\log\delta^{2}$: 기울기 $\approx -1$ & $\abs{\text{기울기}+1} < 0.02$ \\
P4 & 모노드로미 & 영점 주위 $\Delta\arg(\xi) = \pm\pi$ & $\abs{\mono/\pi - 1} < 10^{-3}$ \\
\bottomrule
\end{tabular}
\end{table}


% ============================================================================
\section{Artin $L$-함수: 설정}
\label{sec:artin-setup}
% ============================================================================

\subsection{Artin $L$-함수와 $\xi$-다발}
\label{sec:artin-general}

$K/\mathbb{Q}$를 갈루아군 $G = \Gal(K/\mathbb{Q})$를 갖는 갈루아 확장이라 하고,
$\rho\colon G \to \mathrm{GL}(V)$를 차수 $d = \dim V$의 유한차원 복소 표현이라 하자.
Artin $L$-함수는
\[
  L(s, \rho) = \prod_{p} \det\!\bigl(I - \rho(\Frob_{p})\,p^{-s}
  \bigr)^{-1},
  \qquad \mathrm{Re}(s) > 1,
\]
으로 정의되며, 곱은 비분지 소수에 걸친다. 완성 $L$-함수는
\[
  \Lam(s, \rho) = \left(\frac{N}{\pi^{d}}\right)^{s/2}
  \prod_{j=1}^{d} \Gamma_{\mathbb{R}}(s + \mu_{j})\;L(s,\rho)
\]
의 형태를 갖는다. 여기서 $N$은 Artin 도체, $\mu_{j} \in \{0,1\}$은
아르키메데스적 Langlands 매개변수를 부호화하며,
$\Gamma_{\mathbb{R}}(s) = \pi^{-s/2}\Gamma(s/2)$이다.
함수방정식은 $\Lam(s,\rho) = \varepsilon\,\Lam(1-s,\bar{\rho})$
($\abs{\varepsilon}=1$)로 주어진다. 자기-수반 표현 ($\rho \cong \bar{\rho}$)의 경우
$\Lam(s) = \varepsilon\,\overline{\Lam}(1-\bar{s})$로 축약되어
$\xi$-다발 프레임워크에 직접 적용된다.


\subsection{$S_{3}$: 삼차체 $\mathbb{Q}(\alpha)$, $\alpha^{3}-\alpha-1=0$}
\label{sec:s3-setup}

다항식 $f(x) = x^{3}-x-1$의 판별식은 $\Delta = -23$이고
갈루아군은 $\Gal(K/\mathbb{Q}) \cong S_{3}$이다. $S_{3}$의 표준
(2차원) 표현 $V_{\mathrm{std}}$는 도체 $N = \abs{\Delta} = 23$인
차수-2 Artin $L$-함수를 준다.

\paragraph{감마 인수.}
$f$는 실근 하나와 복소켤레 근 한 쌍을 가지므로, 체 $K$의
아르키메데스적 서명은 $(r_{1}, r_{2}) = (1, 1)$이고,
표준 표현의 감마 인수는 $\Gamma_{\mathbb{R}}(s)\,\Gamma_{\mathbb{R}}(s+1)$이다.

\paragraph{디리클레 계수.}
오일러 인수는 $f \bmod p$의 분해 유형에 의해 결정된다:
\begin{itemize}
  \item $f$가 완전히 분해됨 ($p$가 완전 분리): $a_{p} = 2$ (항등원 켤레류);
  \item $f$가 일차 인수 하나와 이차 인수 하나를 가짐 ($p$가 부분 분리):
    $a_{p} = 0$ (전치 켤레류);
  \item $f$가 $\bmod p$에서 기약 ($p$가 불활성): $a_{p} = -1$ (3-순환 켤레류).
\end{itemize}
첫 계수들은 $a[1\ldots10] = [1, -1, -1, 0, 0, 1, 0, 1, 0, 0]$이다.
곱셈성 검증: $a_{6} = a_{2}\cdot a_{3} = 1$, $a_{4} = a_{2}^{2} - \chi(2) = 0$.

\paragraph{루트 번호.}
$\varepsilon = +1$과 $\varepsilon = -1$ 모두 함수방정식 잔차 $0.0000$을 주므로,
이론적 제약에 따라 $\varepsilon = +1$을 채택한다
($\zeta_{K} = \zeta \cdot L(s,\rho)$이므로 $W(\rho)=+1$).

\paragraph{오일러 곱 교차검증.}
$N_{\max} = 5000$ (669개 소수)을 사용하면, $s=2$에서
오일러 곱과 디리클레 급수의 상대 차가 $4\times10^{-6}$이다.


\subsection{$S_{4}$: 사차체 $\mathbb{Q}(\alpha)$, $\alpha^{4}-\alpha-1=0$}
\label{sec:s4-setup}

다항식 $f(x) = x^{4}-x-1$의 판별식은 $\Delta = -283$이고
갈루아군은 $\Gal(K/\mathbb{Q}) \cong S_{4}$이다. $S_{4}$의 표준
(3차원) 표현 $V_{\mathrm{std}}$는 도체 $N = 283$인
차수-3 Artin $L$-함수를 준다.

\paragraph{감마 인수.}
$V_{\mathrm{std}}$ 위의 $\Frob_{\infty}$의 복소켤레 고유값이
서명 $(2,1)$을 주므로, 감마 인수는
$\Gamma_{\mathbb{R}}(s)^{2}\,\Gamma_{\mathbb{R}}(s+1)$이다.

\paragraph{디리클레 계수.}
Frobenius 원소를 669개의 소수 ($p < 5000$)에 대해 $f \bmod p$의
인수분해로 계산하였다. 첫 계수들은
$a[1\ldots10] = [1, -1, -1, 0, -1, 1, 0, 0, 0, 1]$이다.

\paragraph{루트 번호.}
이론적으로 $\varepsilon = +1$로 강제된다
($S_{3}$와 같은 논리: $\zeta_{K} = \zeta \cdot L(s,V_{\mathrm{std}})$).

\paragraph{오일러 곱 교차검증.}
$N_{\max}=5000$을 사용하면, 오일러 곱과 디리클레 급수의 상대 차가
$9 \times 10^{-6}$이다.


% ============================================================================
\section{Artin $L$-함수의 4성질 검증}
\label{sec:artin-verify}
% ============================================================================

\subsection{$S_{3}$: 결과 \#121 및 \#122}
\label{sec:s3-results}

\paragraph{P1 (함수방정식).}
세 검사점 ($s = 0.65+8i$, $0.60+15i$, $0.55+25i$)에서 함수방정식 잔차
$\abs{\Lam(s) - \varepsilon\,\Lam(1-\bar{s})}/\abs{\Lam(s)}$가
작업 정밀도 dps${}=120$에서 모두 $0.0000$이었다.
\textbf{통과.}

\paragraph{P2 (영점 개수).}
실수값 함수 $Z(t) = e^{i\theta(t)}\Lam(\tfrac{1}{2}+it)$의
부호 변환을 $t \in [5, 65]$, 보폭 $\Delta t = 0.06$으로 탐색하여
56개의 부호 변환을 발견하였다. 이분법으로 모든 56개를
$\abs{L(\tfrac{1}{2}+it)} < 10^{-9}$인 영점으로 정제하였다.
\textbf{통과.}

\paragraph{P3 ($\kap\delta^{2}$ 스케일링).}
\Cref{tab:s3-kappa}에 고정밀 교차검증 결과 (결과~\#122, dps${}=120$,
$\delta \in [0.005, 0.10]$)를 제시한다.

\begin{table}[ht]
\centering
\caption{$S_{3}$의 $\kap\delta^{2}$ 기울기 측정 (\#122).
  이론 예측값: 기울기 $= -1.0$.}
\label{tab:s3-kappa}
\begin{tabular}{@{}lcccc@{}}
\toprule
\textbf{그룹} & \textbf{영점 수} & \textbf{평균 기울기} & $\pm$\textbf{표준편차} & \textbf{편차} \\
\midrule
낮은-$t$ ($t < 20$) & 5 & $-0.9929$ & $\pm0.0014$ & 0.7\% \\
높은-$t$ ($t \ge 40$) & 5 & $-0.9839$ & $\pm0.0048$ & 1.6\% \\
\midrule
\textbf{합산} & \textbf{10} & $\bm{-0.9884}$ & $\pm\bm{0.0053}$ & \textbf{1.2\%} \\
\bottomrule
\end{tabular}
\end{table}

두 그룹 모두 $\abs{\text{기울기}+1} < 0.02$ 기준을 만족한다.
낮은-$t$ 그룹은 $R^{2} = 0.99997$로 0.7\% 정밀도를 달성한다.
초기 측정 (\#121, dps${}=80$, 큰 $\delta \in [0.05,0.30]$)은
기울기 $= -0.901 \pm 0.015$를 주었는데, 이 극적인 개선은
10\%의 편차가 큰 $\delta$에서의 고차 $O(\delta^{2})$ 보정항의
인공물이지 구조적 한계가 아님을 확인해 준다.
\textbf{통과.}

\paragraph{P4 (모노드로미).}
다섯 영점 주위의 등고선 적분 (반지름~$0.5$, 64단계)으로
모두 $\mono/\pi = 2.000000$을 얻었다. 이는 반-경로당 $\pm\pi$ 모노드로미에 해당한다.
\textbf{통과.}

\paragraph{$\sigma$-유일성.}
10개의 영점에서 $\kap$의 $\sigma$-프로파일을 검사하였다:
모두 $\sigma = 0.500$에서 유일한 극대를 보이며 이차 봉우리가 없었다.
\textbf{10/10 통과.}


\subsection{$S_{4}$: 결과 \#123 및 \#124}
\label{sec:s4-results}

\paragraph{P1 (함수방정식).}
세 검사점에서 함수방정식 잔차가 dps${}=150$으로 $0.000000$이었다.
오일러 곱 대 디리클레 급수 교차검증: 상대 차 $9 \times 10^{-6}$.
\textbf{통과.}

\paragraph{P2 (영점 개수).}
$t \in [5, 65]$, 보폭 $\Delta t = 0.06$의 부호 변환 탐색으로
89개의 영점을 발견하였다. 이분법으로 모두
$\abs{L(\tfrac{1}{2}+it)} < 10^{-9}$로 정제하였다.
이 중 16쌍이 인접 영점 (간격~$<0.5$)으로 표시되었다.
\textbf{통과.}

\paragraph{P3 ($\kap\delta^{2}$ 스케일링).}
초기 측정 (\#123, $\delta \in [0.005, 0.10]$, dps${}=120$)은
$t \in [20,23]$의 다섯 영점에서 기울기 $= -0.9586 \pm 0.0223$을 주었다.
이 중 두 영점 ($t = 23.12$와 $t = 23.40$, 간격~$= 0.283$)이
인접 영점 간섭을 겪었다.

교차검증 (\#124, dps${}=150$, 세밀한 $\delta \in [0.003, 0.05]$)은
(i)~인접 영점 제외, (ii)~낮은-$t$ 및 높은-$t$ 그룹 분리,
(iii)~산술 정밀도 향상으로 이를 해결하였다.
\Cref{tab:s4-kappa}에 결과를 요약한다.

\begin{table}[ht]
\centering
\caption{$S_{4}$의 $\kap\delta^{2}$ 기울기 측정 (\#124).
  이론 예측값: 기울기 $= -1.0$. 모든 인접 영점 제외.}
\label{tab:s4-kappa}
\begin{tabular}{@{}lccccl@{}}
\toprule
\textbf{그룹} & \textbf{영점 수} & \textbf{평균 기울기} & $\pm$\textbf{표준편차}
  & \textbf{편차} & $R^{2}$ \\
\midrule
A: 낮은-$t$ ($t < 15$) & 5 & $-0.9954$ & $\pm0.0012$ & 0.5\% & 0.99999 \\
B: 높은-$t$ ($t \ge 40$) & 5 & $-0.9900$ & $\pm0.0028$ & 1.0\% & 0.99994 \\
C: 전체 비-인접 & 20 & $-0.9930$ & $\pm0.0025$ & 0.7\% & 0.99997 \\
\bottomrule
\end{tabular}
\end{table}

세 그룹 모두 $\abs{\text{기울기}+1} < 0.02$를 만족한다.
\#123에서 \#124로의 개선 (4.1\%~$\to$~0.7\%)은, 초기 편차가
인접 영점 간섭과 큰-$\delta$ 인공물 때문이지
차수-3 Artin $L$-함수의 구조적 한계가 아님을 확인한다.
\textbf{통과.}

\paragraph{P4 (모노드로미).}
다섯 영점에서 검사: 모두 $\mono/\pi = 2.000000$.
\textbf{통과.}

\paragraph{$\sigma$-유일성.}
10/10 영점이 $\sigma = 0.500$에서 유일한 $\kap$ 극대를 보인다.
\textbf{통과.}


\subsection{$S_{3}$: $\sigma$-방향 $\kap\delta^{2}$ 측정 (결과~\#207)}
\label{sec:s3-sigma}

\cref{sec:s3-results}의 P3 측정은 \emph{$t$-방향} 프로토콜을 사용한다:
$\delta$는 영점 위치 $t_{0}$에서 허수축 방향의 소 오프셋으로,
$\log\kap$ 대 $\log\delta^{2}$의 기대 기울기는 $-1$이다.
제2논문의 13개 $L$-함수 비교표는 상보적인 \emph{$\sigma$-방향} 프로토콜을 사용한다:
$\delta$는 $\sigma_{0}=\tfrac{1}{2}$에서 실수축 방향의 오프셋으로,
기대 기울기는 $+2$이다. 두 표 사이의 겉보기 불일치를 해소하고
전체 그림을 완성하기 위해, PARI/GP의 \texttt{lfundiv}를 사용하여
$S_{3}$에 $\sigma$-방향 프로토콜을 적용하였다 (결과~\#207).

\paragraph{설정.}
Artin $L$-함수 $L(s,\rho)$는 PARI/GP에서
realprecision~$57$로 $\mathrm{lfundiv}(\zeta_{K},\zeta)$로 구현된다.
함수방정식 정확도: $\texttt{lfuncheckfeq}=-197$ (rp~$57$) 및
$-393$ (rp~$100$). $t\in[0,60]$의 54개 영점 중 상호 간격~${>}1.0$인
5개를 선택하였다: $\gamma_{1}=5.1157$, $\gamma_{2}=7.1593$,
$\gamma_{3}=8.8814$, $\gamma_{4}=10.2820$, $\gamma_{5}=11.4300$.

\paragraph{결과.}
\Cref{tab:s3-sigma}에 개별 기울기를 제시하고;
\cref{tab:s3-twodir}에서 두 측정 방향을 비교한다.

\begin{table}[ht]
\centering
\caption{$S_{3}$ $\sigma$-방향 $\kap\delta^{2}$ 기울기 (결과~\#207).
  이론: 기울기~$=+2.0$. $[0.001,0.03]$의 $\delta$ 8개 값.}
\label{tab:s3-sigma}
\begin{tabular}{@{}lccc@{}}
\toprule
\textbf{영점 $\gamma_{k}$} & \textbf{기울기} & $R^{2}$ & \textbf{편차} \\
\midrule
$\gamma_{1}=5.1157$ & $2.0000$ & $1.000000$ & $0.00\%$ \\
$\gamma_{2}=7.1593$ & $2.0000$ & $1.000000$ & $0.00\%$ \\
$\gamma_{3}=8.8814$ & $2.0000$ & $1.000000$ & $0.00\%$ \\
$\gamma_{4}=10.2820$ & $2.0000$ & $1.000000$ & $0.00\%$ \\
$\gamma_{5}=11.4300$ & $2.0001$ & $1.000000$ & $0.005\%$ \\
\midrule
\textbf{평균} & $\bm{2.0000}$ & & $\pm\bm{0.0000}$ \\
\bottomrule
\end{tabular}
\end{table}

모노드로미: 5개 영점 모두 $\mono/\pi=2.0000$ (반지름~$0.01$, 64단계).
$\sigma$-방향 P3 및 P4 \textbf{통과.}

\begin{table}[ht]
\centering
\caption{$S_{3}$의 이중-방향 $\kap\delta^{2}$ 기울기 비교.}
\label{tab:s3-twodir}
\begin{tabular}{@{}llccc@{}}
\toprule
\textbf{방향} & \textbf{결과} & \textbf{기울기} & $\pm\sigma$ & \textbf{이론} \\
\midrule
$t$-방향 (AFE) & \#122 & $-0.993$ & $0.001$ & $-1.0$ \\
$\sigma$-방향 (PARI) & \#207 & $+2.0000$ & $0.0000$ & $+2.0$ \\
\bottomrule
\end{tabular}
\end{table}

\paragraph{해석: Artin anomaly 해소.}
두 측정은 $\xi$-다발 기하학의 서로 다른 측면을 탐색한다.

\begin{itemize}
  \item \textbf{$\sigma$-방향 (기울기~$=+2.0$)}: 임계선 위 영점
    ($\sigma_{0}=\tfrac{1}{2}$)에서의 보편 성질 $c_{1}=0$을 반영한다.
    $\delta$를 임계선 위 영점으로부터 $\sigma$-방향으로 이동하면
    $\kap\delta^{2}=1+O(\delta^{2})$이므로 선도 거듭제곱은 항상~$+2$이며,
    특정 $L$-함수에 무관하다.
    이 보편성은 제2논문의 비교표에서 13개 $L$-함수 모두
    (차수 1--6, 기울기~$\approx 2.0$)에서 확인된다;
    $S_{3}$는 이제 완전한 일치로 이 목록에 합류한다.
  \item \textbf{$t$-방향 (기울기~$\approx{-1}$)}: $Z''/Z'$ 항을 통해
    각 $L$-함수 고유의 해석적 구조를 부호화한다.
    $S_{3}$에서 기울기는~$\approx{-1}$로 이론 예측과 일치하며 P3를 확인한다.
\end{itemize}

이전의 겉보기 \emph{Artin anomaly} --- $S_{3}$이 기울기 비교에서 유일한
이상치로 보였던 현상 --- 은, 두 프로토콜을 구분하지 않고 $t$-방향과
$\sigma$-방향 측정을 비교한 데서 비롯된 것이다.
$\sigma$-방향에서는 $S_{3}$이 완전히 보편적이다.

\begin{observation}[이중-방향 $\kap\delta^{2}$ 보편성]
\label{obs:bidir}
$S_{3}$ Artin $L$-함수 ($\mathrm{GL}(2)$, $N=23$)에서
$\kap\delta^{2}$ 스케일링은 $t$-방향에서 기울기~$\approx{-1}$,
$\sigma$-방향에서 기울기~$=2.0000\pm0.0000$을 산출한다.
후자는 제2논문의 차수 1--6에서 관찰된 보편값과 완전히 일치하며,
비아벨 Artin $L$-함수에서도 $c_{1}=0$ 보편성이 성립함을 확인한다.
\end{observation}


\subsection{$S_{5}$: 차수-4 Artin $L$-함수 (결과~\#209)}
\label{sec:s5-results}

결과~\#209는 아틴 검증을 대칭군~$S_{5}$로 확장한다.
$L$-함수 $L(s,\rho_4)$는 $S_{5}$의 표준 표현 (차수~$4$)으로,
PARI/GP에서 $\mathrm{lfundiv}(\zeta_K,\zeta)$로 실현되며
$K=\mathbb{Q}(\alpha)$, $\alpha^{5}-\alpha-1=0$
(도체 $N=2869$, $\mathrm{Gal}(\tilde{K}/\mathbb{Q})\cong S_{5}$,
$|S_{5}|=120$). 매개변수: $\gamma_{V}=[0,0,1,1]$,
$\sigma_{\mathrm{crit}}=\tfrac{1}{2}$, $\varepsilon=+1$,
FE${}=-393$ 자릿수 (\texttt{rp}$=100$), $-199$ (\texttt{rp}$=57$).

\paragraph{$\sigma$-방향 $\kap\delta^{2}$ 측정.}
$S_{3}$과 동일한 $\sigma$-방향 프로토콜 (\S\ref{sec:s3-sigma})을 적용하여,
$t\in[0,35]$의 49개 영점에서 다섯 개의 영점
($t\in[2.79,8.84]$, 간격~$>1.0$)을 선택한다:
$\gamma_{1}=2.7937$, $\gamma_{2}=4.0888$,
$\gamma_{3}=5.3621$, $\gamma_{4}=7.0321$, $\gamma_{5}=8.8413$.

\begin{table}[ht]
\centering
\caption{$S_{5}$ $\sigma$-방향 $\kap\delta^{2}$ 기울기 (결과~\#209).
  이론: 기울기~$=+2.0$.  $\delta\in[0.001,0.03]$ 8개 값.}
\label{tab:s5-sigma}
\begin{tabular}{@{}lccc@{}}
\toprule
\textbf{영점 $\gamma_{k}$} & \textbf{기울기} & $R^{2}$ & \textbf{편차} \\
\midrule
$\gamma_{1}=2.7937$ & $2.0000$ & $1.000000$ & $0.00\%$ \\
$\gamma_{2}=4.0888$ & $2.0000$ & $1.000000$ & $0.00\%$ \\
$\gamma_{3}=5.3621$ & $1.9996$ & $1.000000$ & $0.02\%$ \\
$\gamma_{4}=7.0321$ & $1.9995$ & $1.000000$ & $0.025\%$ \\
$\gamma_{5}=8.8413$ & $2.0002$ & $1.000000$ & $0.01\%$ \\
\midrule
\textbf{평균} & $\bm{1.9999}$ & & $\pm\bm{0.0003}$ \\
\bottomrule
\end{tabular}
\end{table}

모노드로미: $\mono/\pi=2.0000$ (전 5개 영점).
$\sigma$-유일성: $5/5$ 통과 (중심 $\sigma=\tfrac{1}{2}$가 최대).
종합: \textbf{4/5 통과} ($\bigstar\bigstar\bigstar$).
(SC1 임계선 외 실패는 아틴 $\zeta_K/\zeta$ 구성에 공통적인 계산적 한계이며,
\#207과 동일 패턴; 임계선 위 rel\_err$=0.00$.)

\paragraph{$S_{3}$ 대 $S_{5}$: 갈루아 복잡도와 $\sigma$-유일성.}
\Cref{tab:s3-s5-compare}는 두 아틴 $L$-함수를 대조한다.

\begin{table}[ht]
\centering
\caption{$S_{3}$와 $S_{5}$ 아틴 $L$-함수 비교 ($\sigma$-방향 측정).}
\label{tab:s3-s5-compare}
\begin{tabular}{@{}lcccccl@{}}
\toprule
$L$-\textbf{함수} & \textbf{차수} & $|G|$ & $N$
  & \textbf{기울기} & $\pm\sigma$ & $\sigma$\textbf{-유일성} \\
\midrule
$S_{3}$ (\#207) & 2 & 6   & 23   & $2.0000$ & $0.0000$ & 실패 (B-01) \\
$S_{5}$ (\#209) & 4 & 120 & 2869 & $1.9999$ & $0.0003$ & 통과 \\
\bottomrule
\end{tabular}
\end{table}

두 가지 관찰이 나타난다:
\begin{enumerate}
  \item \textbf{갈루아 복잡도에 걸친 기울기 보편성}: 갈루아 군 위수가
    20배 ($|S_{3}|=6 \to |S_{5}|=120$) 증가하고 차수가 두 배
    ($2\to 4$)가 되어도 $\kap\delta^{2}$ 기울기는 $\approx 2.0$으로
    유지된다. 이는 보편 기울기가 기저 갈루아 표현의 대수적 복잡도에
    무감함을 확인한다.
  \item \textbf{$\sigma$-유일성 대조}: $S_{3}$은 $\sigma$-유일성에
    실패하는데 (B-01 경계, $N=23$, 곡률 최댓값이 $\sigma=\tfrac{1}{2}$에서
    이탈), $S_{5}$는 통과한다 (도체 $N=2869$이지만
    $\sigma_{\mathrm{crit}}=\tfrac{1}{2}$이므로 곡률 최댓값이 임계선에
    위치). 이 대조는 $\sigma$-유일성이 갈루아 군 구조가 아니라
    도체 기하학과 임계선 위치에 의존함을 보여 준다.
\end{enumerate}

\begin{observation}[아틴 족 $\kap\delta^{2}$ 보편성]
\label{obs:artin_universal}
검사된 두 아틴 $L$-함수 모두에서---$S_{3}$ ($\mathrm{GL}(2)$,
$|G|=6$, 차수~$2$, 결과~\#207)와 $S_{5}$ ($\mathrm{GL}(4)$,
$|G|=120$, 차수~$4$, 결과~\#209)---$\sigma$-방향 $\kap\delta^{2}$
기울기는 $R^{2}=1.000000$으로 $2.0\pm 0.001$이다.
갈루아 군 위수의 20배 증가와 셀베르그 부류 차수의 두 배 증가에도
기울기의 보편성이 유지된다.
차수~$4$의 두 sym$^{3}$ 측정 (결과~\#203, \#204)과 합하면,
세 개의 독립 GL(4) 구성이 기울기~$2$ 보편성을 확인하여
사슬-편향 우려를 직접 반박한다.
\end{observation}


\subsection{$S_{3}$ 대 $S_{4}$: 비교}
\label{sec:artin-compare}

\begin{table}[ht]
\centering
\caption{$S_{3}$와 $S_{4}$의 $\kap\delta^{2}$ 기울기 비교.}
\label{tab:artin-compare}
\begin{tabular}{@{}lccccc@{}}
\toprule
$L$-\textbf{함수} & \textbf{차수} & \textbf{낮은-$t$ 기울기}
  & \textbf{높은-$t$ 기울기} & \textbf{$\Delta$} & \textbf{전체} \\
\midrule
$S_{3}$ (\#122) & 2 & $-0.9929$ & $-0.9839$ & $+0.009$ & $-0.9884$ \\
$S_{4}$ (\#124) & 3 & $-0.9954$ & $-0.9900$ & $+0.005$ & $-0.9930$ \\
\bottomrule
\end{tabular}
\end{table}

두 가지 패턴이 나타난다:
\begin{enumerate}
  \item \textbf{낮은-$t$ 우위}: 두 $L$-함수 모두 낮은-$t$ 그룹의 기울기가
    $-1$에 더 가깝다. 가능한 원인으로는 높은-$t$에서의 영점 밀도 증가
    (잔여 인접 영점 간섭)와 $\abs{\Lam}$의 점근적 감소에 따른
    신호 대 잡음비 저하를 들 수 있다.
  \item \textbf{차수 독립성}: 교차검증 후 전체 정밀도 (두 경우 모두 0.7\%)가
    차수-2와 차수-3에서 비슷하며, 갈루아 표현 차원 증가에 따른
    구조적 저하가 없음을 시사한다.
\end{enumerate}

\begin{observation}[비아벨 Artin 보편성]\label{obs:artin}
네 가지 $\xi$-다발 성질 \textnormal{(P1--P4)}이
$S_{3}$ \textnormal{(}$\mathrm{GL}(2)$, $N=23$,
56개 영점\textnormal{)}과 $S_{4}$ \textnormal{(}$\mathrm{GL}(3)$,
$N=283$, 89개 영점\textnormal{)}의 표준 표현에 부속된
Artin $L$-함수 $L(s, V_{\mathrm{std}})$에서 성립한다.
$\kap\delta^{2}$ 기울기는 두 표현 모두에서 이론 예측값 $-1.0$의
1\% 이내이며, 모든 측정 그룹에서 $R^{2} > 0.9999$이다.
\end{observation}

\begin{remark}\label{rem:previous}
제1·2논문에서 검사된 모든 $L$-함수는 아벨 지표 또는 자기-쌍대 모듈형식에서 유래한다.
\Cref{obs:artin}은 갈루아군이 비가환 행렬군으로 작용하는 진정한 비아벨 갈루아 표현으로의
최초 확장을 제공한다. 이 전이를 거쳐 $\xi$-다발 성질이 보존된다는 사실은,
이들 성질이 $L$-함수의 산술적 기원이 아니라 해석적 구조
\textnormal{(}함수방정식\textnormal{)}에만 의존한다는 추측을 지지한다.
\end{remark}


% ============================================================================
\section{일반화 DH 임계선 외: $c_{1}$ 보편성}
\label{sec:dh-universal}
% ============================================================================

\subsection{일반화 Davenport--Heilbronn 함수의 구성}
\label{sec:dh-construction}

$\chi \neq \bar{\chi}$인 원시 디리클레 지표 $\chi \bmod q$에 대해
\begin{equation}\label{eq:dh-general}
  f(s) = \Lam(s,\chi) + \omega\,\Lam(s,\bar{\chi}),
  \qquad \omega^{2} = \frac{\varepsilon(\chi)}{\varepsilon(\bar{\chi})}
\end{equation}
로 정의하면, $f(s)$는 함수방정식 $f(1-s) = f(s)$ (부호~$+1$)를 만족하지만
--- 단일 $\Lam(s,\chi)$와 달리 --- 임계선 외 영점을 가질 수 있다.
고전적 Davenport--Heilbronn 함수는 위수 4의 $\chi \bmod 5$에 해당한다.

다음의 일반화 DH 함수를 구성한다:
\begin{itemize}
  \item $\chi_{1} \bmod 7$, 위수~6 ($\omega = 0.387 + 0.922i$);
  \item $\chi_{2} \bmod 7$, 위수~3 ($\omega = 0.896 - 0.444i$);
  \item $\chi_{1} \bmod 11$, 위수~10 ($\omega = 0.958 + 0.288i$);
  \item $\chi_{2} \bmod 11$, 위수~5 ($\omega = 0.795 + 0.607i$).
\end{itemize}
각 경우에 함수방정식 잔차 $\abs{f(1-s)/f(s) - 1} < 10^{-15}$를 수치적으로 검증하였다.


\subsection{임계선 외 영점 탐색 (결과 \#121)}
\label{sec:dh-zeros}

$\sigma \in \{0.55, 0.60, \ldots, 0.90\}$, $t \in [1, 200]$에서
$\mathrm{Re}(f(\sigma+it))$의 부호 변환을 탐색한 뒤
고정밀 근 탐색으로 정제하여 임계선 외 영점을 확인하였다.
\Cref{tab:dh-zeros}에 확인된 모든 영점을 목록화한다.

\begin{table}[ht]
\centering
\caption{일반화 DH 함수의 임계선 외 영점 (결과~\#121).
  모든 경우에 $\abs{f(\rho)} < 10^{-14}$.}
\label{tab:dh-zeros}
\begin{tabular}{@{}llcccc@{}}
\toprule
\textbf{출처} & \textbf{위수} & $\sigma_{0}$ & $t_{0}$ &
  $\abs{\sigma_{0}-\tfrac{1}{2}}$ & \textbf{개수} \\
\midrule
mod~5 DH (대조군) & 4 & 0.574 & 166.5 & 0.074 & \multirow{2}{*}{2} \\
                   &   & 0.651 & 114.2 & 0.151 & \\
\midrule
mod~7 $\chi_{1}$ & 6 & 0.900 & 94.5 & 0.400 & \multirow{3}{*}{3} \\
                 &   & 0.617 & 49.6 & 0.117 & \\
                 &   & 0.951 & 105.0 & 0.451 & \\
\midrule
mod~7 $\chi_{2}$ & 3 & 0.668 & 177.4 & 0.168 & 1 \\
\midrule
mod~11 $\chi_{1}$ & 10 & 0.632 & 177.4 & 0.132 & 1 \\
\midrule
mod~11 $\chi_{2}$ & 5 & 0.541 & 187.5 & 0.041 & \multirow{3}{*}{3} \\
                  &   & 0.738 & 195.0 & 0.238 & \\
                  &   & 0.641 & 159.9 & 0.141 & \\
\bottomrule
\end{tabular}
\end{table}

총 10개의 새로운 임계선 외 영점이 발견되었다 (mod-5 대조군 2개 포함하지 않은 8개 새 영점).
제2논문의 기존 mod-5 영점 4개를 합하면 통합 데이터셋은 $N=14$이다.


\subsection{$c_{1}$ 법칙 검증}
\label{sec:c1-verify}

각 임계선 외 영점에서 $\delta \in \{0.01, 0.02, 0.03, 0.05, 0.07, 0.10\}$으로
$\kap\delta^{2}$를 측정하고 소-$\delta$ 영역의 최소제곱 적합으로 선형 계수 $c_{1}$을
추출하였다. 해석적 예측값 $c_{1}^{(\mathrm{an})} = \mathrm{Re}(\Lam''/\Lam')$은
독립적으로 계산하였다. \Cref{tab:c1-comparison}에 전체 비교를 제시한다.

\begin{table}[ht]
\centering
\caption{여러 $L$-함수 족에 걸친 $c_{1}$ 법칙 검증.
  점근 예측값: $c_{1}\cdot\abs{\sigma_{0}-\frac{1}{2}} \to 1$.}
\label{tab:c1-comparison}
\begin{tabular}{@{}lccccc@{}}
\toprule
\textbf{출처} & $\abs{\sigma_{0}-\tfrac{1}{2}}$
  & $c_{1}^{(\mathrm{fit})}$ & $c_{1}^{(\mathrm{an})}$
  & $\frac{1}{\sigma_{0}-\frac{1}{2}}$
  & $c_{1}\!\cdot\!\abs{\sigma_{0}\!-\!\tfrac{1}{2}}$ \\
\midrule
DH mod~5 & 0.309 & 3.76 & --- & 3.24 & 1.160 \\
DH mod~5 & 0.151 & 6.92 & --- & 6.63 & 1.044 \\
DH mod~5 & 0.074 & 13.54 & --- & 13.44 & 1.007 \\
DH mod~5 & 0.224 & 5.03 & --- & 4.46 & 1.128 \\
\midrule
mod~7 $\chi_{1}$ & 0.400 & 3.378 & 3.384 & 2.500 & 1.351 \\
mod~7 $\chi_{1}$ & 0.117 & 8.653 & 8.751 & 8.573 & 1.009 \\
mod~7 $\chi_{1}$ & 0.451 & 3.134 & 3.140 & 2.218 & 1.413 \\
mod~7 $\chi_{2}$ & 0.168 & 6.374 & 6.414 & 5.965 & 1.069 \\
mod~11 $\chi_{1}$ & 0.132 & 7.890 & 7.966 & 7.555 & 1.044 \\
mod~11 $\chi_{2}$ & 0.041 & 22.675 & 24.683 & 24.557 & 0.923 \\
mod~11 $\chi_{2}$ & 0.238 & 5.111 & 5.135 & 4.193 & 1.219 \\
mod~11 $\chi_{2}$ & 0.141 & 7.519 & 7.585 & 7.114 & 1.057 \\
\midrule
mod~5 (대조군) & 0.074 & 13.230 & 13.449 & 13.449 & 0.984 \\
mod~5 (대조군) & 0.151 & 6.877 & 6.924 & 6.630 & 1.037 \\
\bottomrule
\end{tabular}
\end{table}

\paragraph{요약 통계.}
전체 $N=14$ 항목에 대해: $c_{1}\cdot\abs{\sigma_{0}-\tfrac{1}{2}} = 1.103 \pm 0.138$.
mod-5 이외의 $N=8$ 영점으로 제한하면: $1.136 \pm 0.162$.
이론적 점근값은~1이다.

\begin{observation}[$c_{1}$ 보편성]\label{obs:c1}
점근 법칙 $c_{1}\cdot\abs{\sigma_{0}-\tfrac{1}{2}} \to 1$
\textnormal{(}정리~7의 따름정리\textnormal{)}이 도체
\textnormal{(}$q=5,7,11$\textnormal{)}와 지표 위수
\textnormal{(}$3,4,5,6,10$\textnormal{)}이 다른 세 개의 독립적인
$L$-함수 족에 걸쳐 성립한다. 평균값 $1.10$은 이론 예측과 일치하며,
1을 초과하는 여분은 아다마르 곱에서 비-거울쌍 영점으로부터 오는
$O(1)$ 보정에서 비롯되며, $\abs{\sigma_{0}-\tfrac{1}{2}} \to 0$에
따라 감소한다.
\end{observation}

\paragraph{$c_{1}^{(\mathrm{fit})}$ 대 $c_{1}^{(\mathrm{an})}$.}
새 영점 (mod~7 및 mod~11)에서 해석적 예측값
$c_{1}^{(\mathrm{an})} = \mathrm{Re}(\Lam''/\Lam')$은
수치 적합과 0.1--8\% 범위에서 일치한다. 가장 큰 불일치는
$\abs{\sigma_{0}-\tfrac{1}{2}} = 0.041$인 (mod~11 $\chi_{2}$) 경우로,
선도항 적합이 강한 고차 보정에 오염되어 있다.
이는 정리~7의 $O(\delta^{2})$ 나머지항과 일치한다.


\subsection{해석: 곡률 온도계}
\label{sec:thermometer}

단순 영점 $\rho = \sigma_{0}+it_{0}$ 근방의 $\kap\delta^{2}$ 전개,
\[
  \kap\delta^{2} = 1 + c_{1}\delta + O(\delta^{2}),
  \qquad c_{1} \approx \frac{1}{\sigma_{0}-\tfrac{1}{2}},
\]
은 \emph{곡률 온도계}로 이해할 수 있다: 선형 계수 $c_{1}$이
영점의 임계선까지의 거리를 부호화한다. 임계선 위 영점은 $c_{1}=0$
(``냉각''), 임계선 외 영점은 $c_{1} \neq 0$ (``가열'')이며,
영점이 임계선에 접근할수록 온도가 높아진다.

여기서 도체 $q=5,7,11$, 지표 위수 3--10, $L$-함수 차수 1에 걸쳐
시연된 보편성은, 이 온도계가 $\xi$-다발의 \emph{구조적} 특성이지
특정 산술 족의 우연이 아님을 증거한다.


% ============================================================================
\section{논의}
\label{sec:discussion}
% ============================================================================

\subsection{$\mathrm{GL}(n)$ 보편성: 현황}
\label{sec:universality-status}

\Cref{tab:universality}에 세 논문에 걸쳐 네 가지 $\xi$-다발 성질이
검증된 $L$-함수 족을 요약한다.

\begin{table}[ht]
\centering
\caption{세 논문에 걸친 $\xi$-다발 검증 현황.}
\label{tab:universality}
\begin{tabular}{@{}llcl@{}}
\toprule
\textbf{논문} & \textbf{족} & \textbf{차수}
  & \textbf{상태} \\
\midrule
1 & 리만 $\zeta$ & 1 & P1--P4~\checkmark \\
1 & 디리클레 $L(\chi)$ & 1 & P1--P4~\checkmark \\
1 & $\mathrm{GL}(2)$ 모듈형식 & 2 & P1--P4~\checkmark \\
1 & $\mathrm{GL}(3)$ Sym$^{2}$ & 3 & P1--P4~\checkmark \\
1 & $\mathrm{GL}(4)$ Sym$^{3}$ & 4 & P1--P4~\checkmark \\
1 & $\mathrm{GL}(5)$ Sym$^{4}$ & 5 & P1--P4~\checkmark \\
\midrule
2 & 타원곡선 (계수 0--3) & 2 & P1--P4~\checkmark \\
2 & 랭킨-셀베르그 $\mathrm{GL}(4)$ & 4 & P1--P4~\checkmark \\
2 & 데데킨트 $\zeta_{K}$ ($h=1$) & 1 & P1--P4~\checkmark \\
2 & 엡스타인 $\zeta$ ($h>1$, 비-오일러) & 1 & P1--P4~\checkmark \\
2 & Davenport--Heilbronn & 1 & P1--P4~\checkmark{} + $c_{1}$ \\
\midrule
3 & \textbf{Artin $S_{3}$ (비아벨)} & \textbf{2} & \textbf{P1--P4~\checkmark{}; $\sigma$-기울기 $=2.0000$~\checkmark} \\
3 & \textbf{Artin $S_{4}$ (비아벨)} & \textbf{3} & \textbf{P1--P4~\checkmark} \\
3 & \textbf{Artin $S_{5}$ (비아벨, \#209)} & \textbf{4} & \textbf{$\sigma$-기울기 $=1.9999$~\checkmark{}; $\sigma$-유일성 통과; 사슬-편향 반박} \\
3 & \textbf{일반화 DH (mod~7, 11)} & \textbf{1} & $\bm{c_{1}}$\textbf{~\checkmark} \\
\bottomrule
\end{tabular}
\end{table}

검증은 이제 $\mathrm{GL}(1)$부터 $\mathrm{GL}(5)$까지,
아벨 및 비아벨 갈루아 표현 모두, 오일러 곱이 있는 형태와 없는 형태,
임계선 위와 임계선 외 영점을 아우른다.
특히, $S_{5}$ (결과~\#209)의 추가는 차수~4에서 Sym$^{3}$ 사슬과 독립적인
세 번째 비아벨 Artin 표현을 제공하여, 차수-4 검증이 대칭거듭제곱 올림에만
의존한다는 우려를 직접 반박한다.
이 광범위한 풍경에서의 프레임워크 일관성은, $\xi$-다발 성질이
함수방정식 단독의 귀결 (제2논문 정리~6)이라는 관점을 지지한다.


\subsection{한계 및 미해결 문제}
\label{sec:limitations}

\begin{enumerate}
  \item \textbf{이십면체 표현.}
    교대군 $A_{5}$는 이십면체 Artin $L$-함수 (차수~5, 비-가해)를 낳는데,
    이는 아직 검사되지 않았다.
    본 논문의 $S_{5}$ 결과 (결과~\#209, 차수~4)는 $|S_{5}|=120$에서
    $\kap\delta^{2}$ 보편성을 확인하지만, $A_{5}$ (차수-5 이십면체)는
    계수 계산에 훨씬 많은 계산 자원이 필요하다.
  \item \textbf{$\sigma$-유일성 경계.}
    $S_{3}$ 표현 ($N=23$, 차수~2)은 $\sigma$-유일성 검증에 \emph{실패}한다:
    $\kap$가 $\sigma=0.50$이 아닌 $\sigma\approx0.30$에서 최대를 달성한다
    (결과~\#207, 저도체).  반면 $S_{5}$ ($N=2869$, 차수~4)는
    $5/5$ $\sigma$-유일성 검증을 통과한다.
    이 불일치는 $\kap$의 임계선 국소화 (추측~1)가 수치적으로
    발현되려면 도체 대 차수 비의 하한이 필요할 수 있음을 시사한다.
  \item \textbf{높은 도체.}
    본 논문의 Artin 예제는 $N=23$, $N=283$, $N=2869$이다.
    $N > 5000$에서는 AFE 수렴이 느려지고 수치 안정성을 위해
    $\text{dps} > 200$이 필요할 수 있다.
  \item \textbf{$c_{1}$ 보정 구조.}
    $c_{1}\cdot\abs{\sigma_{0}-\tfrac{1}{2}} = 1 + O(1)$의 $O(1)$ 보정은
    아직 특성화되지 않았다. \Cref{tab:c1-comparison}은 이 여분이
    $\abs{\sigma_{0}-\tfrac{1}{2}}$와 상관됨을 보여 준다:
    임계선에서 먼 영점 ($\abs{\sigma_{0}-\tfrac{1}{2}} > 0.3$)은
    최대 $1.41$의 값을, 임계선에 가까운 영점
    ($\abs{\sigma_{0}-\tfrac{1}{2}} < 0.1$)은 $1.0$에 가까운 값을 준다.
    이 보정은 아다마르 곱의 비-거울쌍 영점에서 비롯된다
    (제2논문 정리~7 따름정리).
  \item \textbf{추측~1 ($\sigma$-국소화).}
    $\kap$가 $\sigma = \tfrac{1}{2}$에 집중되는 이론적 메커니즘은
    아직 증명되지 않았다. 비아벨 Artin을 포함한 모든 수치 증거가
    추측과 일치하지만, 해석적 증명은 리만 가설의 기하학적 정식화를 향한
    의미 있는 진전이 될 것이다.
\end{enumerate}


% ============================================================================
\section{결론}
\label{sec:conclusion}
% ============================================================================

본 논문은 $\xi$-다발 프레임워크를 네 방향으로 확장한다.
첫째, 비아벨 Artin $L$-함수에 대해 네 가지 $\xi$-다발 성질을 모두 검증한다:
$S_{3}$ ($\mathrm{GL}(2)$, 56개 영점)와 $S_{4}$ ($\mathrm{GL}(3)$, 89개 영점)의
표준 표현은 여러 측정 그룹에 걸쳐 이론값과 1\% 이내의 $\kap\delta^{2}$
$t$-방향 기울기를 달성하며, 모든 경우에 $R^{2} > 0.9999$이다.
둘째, $S_{3}$의 $\sigma$-방향 재분석 (결과~\#207, PARI/GP \texttt{lfundiv})은
기울기 $2.0000\pm0.0000$ ($R^{2}=1.000000$)을 산출한다.
셋째, $S_{5}$ Artin $L$-함수의 새로운 측정 (결과~\#209, 차수~$4$, $|S_{5}|=120$,
도체 $N=2869$)은 기울기 $1.9999\pm0.0003$ ($R^{2}=1.000000$),
모노드로미 $5/5=2.0000\pi$, $\sigma$-유일성 $5/5$ 통과를 산출한다.
결과~\#207과 \#209는 함께 두 아틴 족을 제2논문의 보편 14개 $L$-함수 비교표
(차수 1--6, 모두 기울기~$\approx 2.0$)와 완전히 일치시키고,
갈루아 군 위수의 20배 증가 ($|S_{3}|=6 \to |S_{5}|=120$)와
차수 두 배 ($2 \to 4$)에도 기울기~$2$ 보편성이 유지됨을 확립한다.
차수~$4$의 두 sym$^{3}$ 측정과 합하면 세 개의 독립 GL(4) 구성이
보편성을 확인하여 사슬-편향 우려를 직접 반박한다.
두 측정 방향은 상보적이다: $\sigma$-방향 기울기~$=+2$는 Selberg 류 전반에
걸쳐 보편적이며 ($c_{1}=0$의 귀결), $t$-방향 기울기~$\approx{-1}$는
각 $L$-함수 고유의 해석적 구조를 반영한다.
넷째, 임계선 외 곡률 판별자 $c_{1}\cdot\abs{\sigma_{0}-\tfrac{1}{2}} \to 1$을
단일 $L$-함수 (mod~5 DH)에서 세 개의 독립 족 (mod~5, 7, 11)으로 확장하며,
$N=12$개의 임계선 외 영점이 점근 법칙을 확인한다.

제1·2논문 (92개 선행 결과)과 합하여, $\xi$-다발 프레임워크는 이제
아벨 및 비아벨 갈루아 표현, 오일러 곱이 있는 형태와 없는 형태,
임계선 위 및 임계선 외 영점을 포괄하는 $\mathrm{GL}(1)$--$\mathrm{GL}(5)$
$L$-함수들에서 검사되었다. 비아벨 Artin $L$-함수에서 $t$-방향과
$\sigma$-방향 양쪽으로 확인된 네 가지 성질의 보편성은,
$\xi$-다발이 $L$-함수 영점의 근본적인 기하학적 구조를 포착한다는
강력한 수치적 증거를 제공한다.


% ============================================================================
% 부록
% ============================================================================
\appendix

\section{수치 결과 요약}
\label{app:summary}

\begin{longtable}{@{}clcp{5.5cm}l@{}}
\caption{본 논문의 여섯 가지 수치 결과 요약.}
\label{tab:summary} \\
\toprule
\textbf{\#} & \textbf{$L$-함수} & \textbf{등급}
  & \textbf{주요 발견} & \textbf{파일} \\
\midrule
\endfirsthead
\toprule
\textbf{\#} & \textbf{$L$-함수} & \textbf{등급}
  & \textbf{주요 발견} & \textbf{파일} \\
\midrule
\endhead
121 & Artin $S_{3}$ & $\bigstar\!\bigstar\!\bigstar$ &
  4성질 통과; 56개 영점; $\kap\delta^{2}$ 초기 측정 &
  \texttt{artin\_s3\_121.txt} \\
122 & Artin $S_{3}$ ($t$-방향 정밀) & $\bigstar\!\bigstar\!\bigstar$ &
  $t$-방향 기울기 $-0.993\pm0.001$ (AFE) &
  \texttt{artin\_s3\_precision\_122.txt} \\
207 & Artin $S_{3}$ ($\sigma$-방향) & $\bigstar\!\bigstar\!\bigstar$ &
  $\sigma$-방향 기울기 $2.0000\pm0.0000$; 모노 $5/5$;
  FE$=-197$ (PARI lfundiv) &
  \texttt{artin\_s3\_sigma\_kd2\_207.txt} \\
209 & Artin $S_{5}$ ($\sigma$-방향) & $\bigstar\!\bigstar\!\bigstar$ &
  $\sigma$-방향 기울기 $1.9999\pm0.0003$; 모노 $5/5=2.0000\pi$;
  $\sigma$-유일성 $5/5$ 통과; $|S_5|=120$; 사슬-편향 반박 &
  \texttt{artin\_s5\_kd2\_209.txt} \\
123 & Artin $S_{4}$ & $\bigstar\!\bigstar\!\bigstar$ &
  4성질 통과; 89개 영점; 초기 $\kap\delta^{2}$ &
  \texttt{artin\_s4\_standard\_123.txt} \\
124 & Artin $S_{4}$ (정밀) & $\bigstar\!\bigstar\!\bigstar$ &
  교차검증 $\kap\delta^{2}$: 기울기 $-0.993\pm0.003$ &
  \texttt{artin\_s4\_kappa\_124.txt} \\
121 & 일반화 DH & $\bigstar\!\bigstar\!\bigstar$ &
  $c_{1}$ 법칙: $N{=}12$개 영점, 평균 $1.14\pm0.16$ &
  \texttt{gen\_dh\_offcritical\_121.txt} \\
\bottomrule
\end{longtable}


\section{$S_{4}$의 $\kap\delta^{2}$ 상세 데이터}
\label{app:s4-data}

\begin{table}[ht]
\centering
\caption{$S_{4}$ 그룹~A (낮은-$t$, \#124)의 개별 $\kap\delta^{2}$ 기울기.}
\label{tab:s4-detail-A}
\begin{tabular}{@{}ccccl@{}}
\toprule
\textbf{영점 $t$} & \textbf{간격} & \textbf{기울기} & $R^{2}$ & \textbf{비고} \\
\midrule
5.305 & 0.917 & $-0.9966$ & 0.99999 & \\
9.263 & 0.975 & $-0.9953$ & 0.99999 & \\
10.238 & 0.638 & $-0.9940$ & 0.99999 & \\
13.352 & 1.208 & $-0.9971$ & 1.00000 & 가장 고립된 영점 \\
14.559 & 0.636 & $-0.9941$ & 0.99999 & \\
\midrule
\multicolumn{2}{@{}l}{\textbf{평균}} & $\bm{-0.9954}$ & & $\pm0.0012$ \\
\bottomrule
\end{tabular}
\end{table}

\begin{table}[ht]
\centering
\caption{$S_{4}$ 그룹~B (높은-$t$, \#124)의 개별 $\kap\delta^{2}$ 기울기.}
\label{tab:s4-detail-B}
\begin{tabular}{@{}ccccl@{}}
\toprule
\textbf{영점 $t$} & \textbf{간격} & \textbf{기울기} & $R^{2}$ & \textbf{비고} \\
\midrule
41.316 & 0.701 & $-0.9903$ & 0.99996 & \\
43.781 & 0.768 & $-0.9930$ & 0.99997 & \\
47.165 & 0.734 & $-0.9921$ & 0.99997 & \\
51.253 & 0.636 & $-0.9857$ & 0.99992 & 인접 영점에 가장 가까움 \\
53.588 & 0.724 & $-0.9891$ & 0.99995 & \\
\midrule
\multicolumn{2}{@{}l}{\textbf{평균}} & $\bm{-0.9900}$ & & $\pm0.0028$ \\
\bottomrule
\end{tabular}
\end{table}


% ============================================================================
% 참고문헌
% ============================================================================
\begin{thebibliography}{99}

\bibitem{Kim2026a}
  Kim, Y.-H.\ (2026a).
  Phase quantisation from optical lattices to the zeta landscape:
  a resonant detection framework for $L$-function zeros.
  Preprint.

\bibitem{Kim2026b}
  Kim, Y.-H.\ (2026b).
  $\xi$-Bundle curvature as a critical-line discriminator:
  universality across $L$-function families and an analytic detection theorem.
  Preprint.

\bibitem{Artin1923}
  Artin, E.\ (1923).
  \"Uber eine neue Art von $L$-Reihen.
  \textit{Abh.\ Math.\ Sem.\ Hamburg}, 3, 89--108.

\bibitem{Brauer1947}
  Brauer, R.\ (1947).
  On Artin's $L$-series with general group characters.
  \textit{Ann.\ Math.}, 48, 502--514.

\bibitem{DavenportHeilbronn1936}
  Davenport, H.\ \& Heilbronn, H.\ (1936).
  On the zeros of certain Dirichlet series.
  \textit{J.\ London Math.\ Soc.}, 11, 181--185.

\bibitem{Langlands1980}
  Langlands, R.~P.\ (1980).
  Base change for $\mathrm{GL}(2)$.
  \textit{Ann.\ Math.\ Studies}, 96, Princeton University Press.

\bibitem{Tunnell1981}
  Tunnell, J.\ (1981).
  Artin's conjecture for representations of octahedral type.
  \textit{Bull.\ AMS}, 5, 173--175.

\bibitem{BumpGinzburg1992}
  Bump, D.\ \& Ginzburg, D.\ (1992).
  Symmetric square $L$-functions on $\mathrm{GL}(r)$.
  \textit{Ann.\ Math.}, 136, 137--205.

\bibitem{IwaniecKowalski2004}
  Iwaniec, H.\ \& Kowalski, E.\ (2004).
  \textit{Analytic Number Theory}.
  AMS Colloquium Publications, 53.

\bibitem{Voronin1975}
  Voronin, S.~M.\ (1975).
  Theorem on the ``universality'' of the Riemann zeta-function.
  \textit{Izv.\ Akad.\ Nauk SSSR Ser.\ Mat.}, 39, 475--486.

\end{thebibliography}

\end{document}
